% ch-eu-causal -- the causal layer of the transparent regime (imported
% from the handbook's causal chapter; body of "A Causal Layer: Does
% Financing Cause Growth?")

On the decision map of \cref{ch01:sec:map}, this chapter serves the one consuming action that
prediction cannot: capital-allocation post-mortems and targeting (``did our financings work?'',
``where would a round move the needle?''). \Cref{ch:eufund} predicts; managers ask a different
question: \emph{if} a company raises a round, does growth follow --- or do growing companies
simply raise? The fitted $\hat m(x)$ cannot answer this, because the dependence it learns on
deal-history features mixes the effect of financing with the selection of good firms into
financing. This chapter builds the minimal causal layer the transparent regime supports. The
regime is favourable for one reason: the statutory fundamentals panel of \cref{ch:eufund} ---
joined point-in-time to the deals and macro tables of \cref{ch:data} --- gives a \emph{rich,
audited, pre-treatment} covariate history, close to the dossier an investor actually screens
on, so adjustment has a realistic chance. Every estimate ships with cluster-bootstrap
uncertainty, an E-value, and two falsification tests, and the whole pipeline is to be validated
on the synthetic world of \cref{ch02:sec:synthetic}, where the true effect is known by
construction (\cref{ch:results-synth}). Against the
defect ledger of \cref{tab:ch02:defects}, the exposures are those of \cref{ch:eufund} ---
survivor selection in the outcome (handled by the same censoring policy) --- plus one new and
dominant one: confounding by latent quality, which is not in the ledger because it is not a data
defect at all, and which this chapter treats as its central object.

\section{Estimands and assumptions}\label{chca:sec:estimands}

Treatment: $A_{i,y} = \ind{\text{firm } i \text{ closed a PE/VC equity round in year } y}$, from
the deals table. Outcome: $Y_{i,y}$, the $h=2$ winsorized EBITDA arc-growth of
\cref{def:success-label}, measured over $(y, y+h]$ so it is strictly post-treatment. Each unit
carries potential outcomes $Y_{i,y}(1)$ and $Y_{i,y}(0)$: growth with and without a round in
year $y$.

\begin{assumption}[SUTVA]\label{ass:causal-sutva}
(i)~No interference: $Y_{i,y}(a)$ does not depend on other firms' treatments; (ii)~single
version of treatment: $A_{i,y}=1$ denotes one well-defined intervention \citep{imbens2015}.
\end{assumption}

\begin{assumption}[Consistency]\label{ass:causal-consistency}
$Y_{i,y} = A_{i,y} Y_{i,y}(1) + (1 - A_{i,y})\, Y_{i,y}(0)$ \citep{hernan2020}.
\end{assumption}

Both parts of \cref{ass:causal-sutva} are approximations here: a round can hurt a direct
competitor's growth (interference; we flag narrow niches), and rounds differ in size and
investor (versions; we deliberately coarsen to ``raised institutional equity'' and treat
dose--response as an extension, not the estimand).

The estimands, in decreasing ambition of the conditioning:
\begin{equation}\label{eq:causal-estimands}
\tau_{\mathrm{ATE}} = \E[Y(1) - Y(0)], \quad
\tau_{\mathrm{ATT}} = \E[Y(1) - Y(0) \st A = 1], \quad
\tau(x) = \E[Y(1) - Y(0) \st X = x].
\end{equation}
For the portfolio question (``did our financings work?'') the ATT is the object; for targeting
(``where would a round move the needle?'') the CATE $\tau(x)$; the ATE is the reporting default
because it is the easiest to estimate stably. Identification requires:

\begin{assumption}[Conditional exchangeability and positivity]
\label{ass:causal-ignorability} There exists an observed pre-treatment set $Z$ with $\bigl(Y(0),
Y(1)\bigr) \indep A \st Z$ and $0 < e(z) = \Prob(A = 1 \st Z = z) < 1$ for all $z$ in the
support \citep{imbens2015,hernan2020}.
\end{assumption}

\Cref{ass:causal-ignorability} is not a modelling choice; it is a claim about the world that the
next section argues from a graph, \cref{chca:sec:sens} stress-tests, and \cref{chca:sec:synth}
verifies where the truth is known.

\section{The causal graph}\label{chca:sec:dag}

Five variables. $U$: latent firm quality (management, product, moat) --- unobserved. $X$: the
observed pre-treatment fundamentals history (levels, growth path, margins, leverage, payroll
over the three years up to $y$), plus sector and age. $M$: the macro state as of year $y$
(rates, CPI YoY, unemployment, as-of vintages). $A$: the treatment. $Y$: the outcome. The edges:
$U \to A$ (investors select quality), $U \to Y$ (quality drives growth), $U \to X$ (quality
leaves tracks in the filings), $X \to A$ (investors screen on the observable dossier), $X \to Y$
(fundamentals persist), $M \to A$ (funding cycles), $M \to Y$ (macro drives growth), and the
causal edge of interest $A \to Y$. \Cref{chca:tab:paths} enumerates every $A$--$Y$ path.

\begin{table}[htbp]
\centering
\caption{Paths between $A$ and $Y$ in the working graph.}
\label{chca:tab:paths}
\begin{tabular}{@{}llll@{}}
\toprule
Path & Type & Status after conditioning on $(X, M)$ \\
\midrule
$A \to Y$ & causal & kept (the estimand) \\
$A \leftarrow X \to Y$ & backdoor & blocked by $X$ \\
$A \leftarrow M \to Y$ & backdoor & blocked by $M$ (sector-year
cells) \\
$A \leftarrow U \to Y$ & backdoor & \emph{not} blockable directly;
attenuated \\
 & & via the proxy path $U \to X$ \\
$A \leftarrow U \to X \to Y$ & backdoor & blocked by $X$ \\
\bottomrule
\end{tabular}
\end{table}

The whole argument of the transparent-regime track sits in the fourth row. $U$ is unobserved, so
the path $A \leftarrow U \to Y$ cannot be blocked exactly. But investors do not observe $U$
either: they observe the dossier, which is essentially $X$ --- audited statements, growth
trajectory, margins, leverage. To the extent that the selection $U \to A$ flows \emph{through}
what is observable about quality, conditioning on a rich $X$ intercepts most of it, and the
residual confounding is the part of quality that moves investors without leaving any track in
three years of filings. In the transparent regime that residual is plausibly small and,
crucially, \emph{boundable} (\cref{chca:sec:sens}). In the opaque regime there is no
fundamentals panel, $X$ collapses to deal marks and coarse firmographics, and the argument
fails --- which is why \cref{ch:proxies} and everything after it change the estimand rather
than pretend the backdoor is closed.

\section{Identification: the backdoor criterion}\label{chca:sec:backdoor}

\begin{definition}[Backdoor criterion, \citealp{pearl1995}]
\label{def:causal-backdoor} A set of observed variables $Z$ satisfies the backdoor criterion
relative to $(A, Y)$ if (i)~no element of $Z$ is a descendant of $A$, and (ii)~$Z$ blocks every
path between $A$ and $Y$ that contains an arrow into $A$.
\end{definition}

\begin{theorem}[Adjustment formula, \citealp{pearl1995}]
\label{thm:causal-adjustment} If $Z$ satisfies the backdoor criterion relative to $(A, Y)$, then
\begin{equation}\label{eq:causal-adjustment}
\E[Y \st \Do(A = a)] = \E_Z\bigl[\, \E[Y \st A = a, Z] \,\bigr],
\end{equation}
the outer expectation over the marginal law of $Z$, and consequently $\tau_{\mathrm{ATE}} =
\E_Z[\E[Y \st A=1, Z] - \E[Y \st A=0, Z]]$.
\end{theorem}

The proof --- via the truncated factorization, an interventions-do-not-move-non-descendants
lemma, and the invariance half of the do-calculus --- is given in \cref{app:proofs}, together
with the three-line potential-outcomes special case under \cref{ass:causal-ignorability} that
this chapter actually operates under.

In the working graph, $Z = (X, M)$ satisfies the criterion \emph{up to} the residual $U$-path
discussed above: formally we proceed under \cref{ass:causal-ignorability} with $Z = (X$ history
over 3 years, sector-year cell$)$ --- the cells absorb $M$ and sector-level funding cycles ---
and quantify the violation rather than assume it away.

\begin{warning}[Bad controls]\label{warn:causal-badcontrols}
More conditioning is not more safety. Three failure modes, all present in this dataset.
(i)~\emph{Mediators}: post-treatment employee counts sit on the path $A \to \text{hiring} \to
Y$; conditioning on them subtracts part of the effect being estimated.
(ii)~\emph{Colliders/descendants}: a subsequent valuation mark is caused both by the round ($A$)
and by quality ($U$); conditioning on it opens the path $A \to \text{mark} \leftarrow U \to Y$
and manufactures spurious association. (iii)~\emph{M-bias}: even a \emph{pre-treatment} variable
can hurt --- if some early covariate is a common effect of one latent cause that drives
selection and another latent cause that drives growth, then conditioning on it opens a
previously closed path between those latents. The rule is graphical, not chronological: the
adjustment set is chosen by \cref{def:causal-backdoor}, never by predictive power, and nothing
dated after the round enters $Z$.
\end{warning}

\section{Estimators, in increasing robustness}\label{sec:causal-estimators}

Throughout, nuisance functions are fitted with the boosted-tree machinery of \cref{ch:eufund}
(same features, same era discipline), written $\mu_a(z) = \E[Y \st A = a, Z = z]$ and $e(z) =
\Prob(A = 1 \st Z = z)$.

\paragraph{(1) Outcome regression (g-formula).} Fit $\hat\mu_1, \hat\mu_0$ and average:
$\hat\tau_{\mathrm{OR}} = n^{-1} \sum_{i,y} \{\hat\mu_1(Z_{i,y}) - \hat\mu_0(Z_{i,y})\}$. Two
flavours: the T-learner (separate models per arm; honest but noisy in the small treated arm) and
the S-learner ($A$ as one feature among $p$; stable but regularization can shrink the
$A$-splits, biasing $\hat\tau$ toward zero). Both inherit every extrapolation weakness of trees
precisely in low-overlap regions, where $\hat\mu_{1}$ is evaluated at control-like $z$ and vice
versa.

\paragraph{(2) Inverse propensity weighting.} Fit $\hat e$, then the Hajek form
\[
\hat\tau_{\mathrm{IPW}} = \frac{\sum_i A_i Y_i w_i^1}{\sum_i A_i w_i^1} - \frac{\sum_i (1-A_i)
Y_i w_i^0}{\sum_i (1-A_i) w_i^0}, \qquad w^1 = \frac{1}{\hat e(Z)}, \quad w^0 = \frac{1}{1 -
\hat e(Z)}.
\]
Mandatory diagnostics before any estimate is read: the histogram of $\hat e$ by treatment arm
(overlap plot); trimming to $\hat e \in [0.02, 0.98]$ with the trimmed fraction reported --- if
it exceeds a few percent, the estimand has silently changed to a trimmed-population effect and
we say so; the effective sample size $\mathrm{ESS} = (\sum w)^2 / \sum w^2$ per arm. IPW uses
the treatment model only, so it fails independently of (1) --- but its variance explodes near
the trimming boundary.

\paragraph{(3) AIPW with cross-fitting (the default).} The doubly robust score combines both
nuisances \citep{chernozhukov2018}:
\begin{equation}\label{eq:causal-aipw}
\psi(O; \mu, e) = \mu_1(Z) - \mu_0(Z)
+ \frac{A\,(Y - \mu_1(Z))}{e(Z)}
- \frac{(1-A)\,(Y - \mu_0(Z))}{1 - e(Z)},
\qquad
\hat\tau_{\mathrm{AIPW}} = \frac{1}{n} \sum \psi_i .
\end{equation}

\begin{proposition}[Double robustness]\label{prop:causal-dr}
Under \cref{ass:causal-consistency,ass:causal-ignorability}, $\E[\psi(O; \mu, e)] =
\tau_{\mathrm{ATE}}$ if \emph{either} $\mu = (\mu_1^*, \mu_0^*)$ is the true outcome function
\emph{or} $e = e^*$ is the true propensity (with the other nuisance arbitrary, subject to
positivity).
\end{proposition}

\begin{proof}
Consider the $a=1$ half, $\psi_1 = \mu_1(Z) + A(Y - \mu_1(Z))/e(Z)$. If $\mu_1 = \mu_1^*$:
$\E[A(Y - \mu_1^*(Z)) \st Z] = e^*(Z)\,\E[Y - \mu_1^*(Z) \st A=1, Z] = 0$, so $\E[\psi_1] =
\E[\mu_1^*(Z)] = \E[Y(1)]$ by \cref{ass:causal-ignorability}. If $e = e^*$: $\E[\psi_1 \st Z] =
\mu_1(Z) + \{e^*(Z)\,\E[Y \st A=1,Z] - e^*(Z)\mu_1(Z)\}/e^*(Z) = \E[Y \st A=1, Z] = \E[Y(1) \st
Z]$, and the arbitrary $\mu_1$ cancels. The $a=0$ half is symmetric; subtract. (Integrability
details, and the influence-function variance that makes AIPW efficient, are in
\cref{app:proofs}.)
\end{proof}

\begin{remark}[Rate condition]\label{rem:causal-rates}
Beyond consistency, valid $\sqrt{n}$ inference informally requires the \emph{product} of
nuisance errors to vanish fast: $\norm{\hat\mu_a - \mu_a^*}_{2} \cdot \norm{\hat e - e^*}_{2} =
o_P(n^{-1/2})$ \citep{chernozhukov2018}. Each nuisance may converge slowly (as boosted trees do)
as long as the product is fast --- this, plus the removal of own-observation overfitting bias by
cross-fitting, is why (3) is the default and not merely a hedge.
\end{remark}

\paragraph{Algorithm (cross-fitted AIPW).}
\begin{enumerate}
\item Partition \emph{firms} into $K = 5$ folds; a firm's rows never straddle folds (within-firm
correlation would otherwise leak nuisance fits into own scores).
\item For each fold $k$: fit $\hat\mu_1^{(-k)}, \hat\mu_0^{(-k)}, \hat e^{(-k)}$ on the other
$K-1$ folds, with the era-aware tuning of \cref{ch:eufund} frozen once, outside the loop.
\item Trim: drop rows of fold $k$ with $\hat e^{(-k)} \notin [0.02, 0.98]$; record the trimmed
fraction.
\item Compute $\psi_{i,y}$ for the surviving rows of fold $k$ via \cref{eq:causal-aipw} using
only out-of-fold nuisances.
\item Pool: $\hat\tau = $ mean of all $\psi_{i,y}$; variance by \cref{sec:causal-inference};
report alongside the overlap plot and trimmed fraction.
\end{enumerate}

\section{Inference under firm clustering}\label{sec:causal-inference}

Rows recur within firms across years, and the scores $\psi_{i,y}$ of a firm are correlated
through persistent firm effects. An i.i.d.\ bootstrap (or the naive variance
$n^{-2}\sum(\psi_{i,y} - \hat\tau)^2$) treats $n$ rows as $n$ independent draws and understates
the variance --- the effective sample size is nearer the number of firms. Two consistent routes;
we run both and report the wider.

\paragraph{(a) Cluster-robust influence-function variance.} Let $S_i = \sum_{y} (\psi_{i,y} -
\hat\tau)$ sum a firm's centred scores. Then
\begin{equation}\label{eq:causal-clustervar}
\widehat{\Var}(\hat\tau) = \frac{1}{n^2} \sum_{i \in \firms}
S_i^{\,2},
\end{equation}
which is the sandwich variance with firms as clusters; normal intervals $\hat\tau \pm
z_{1-\alpha/2}\,\widehat{\Var}^{1/2}$.

\paragraph{(b) Cluster bootstrap \citep{efron1994}.}
\begin{enumerate}
\item For $b = 1, \dots, B$ (200--500): resample firms with replacement; a drawn firm
contributes all its rows.
\item Re-run the \emph{entire} cross-fitting algorithm on the replicate --- nuisance refits
included, tuning still frozen --- to obtain $\hat\tau_b$.
\item Report the percentile interval of $\{\hat\tau_b\}$.
\end{enumerate}
The bootstrap additionally propagates nuisance-estimation noise that \cref{eq:causal-clustervar}
treats as fixed; disagreement between (a) and (b) beyond $\sim$20\% of interval width is itself
a diagnostic of nuisance instability.

\section{Sensitivity and falsification}\label{chca:sec:sens}

\Cref{ass:causal-ignorability} cannot be verified from the data it licenses. It can be
stress-tested. Every reported effect must clear the checklist of \cref{chca:tab:checklist}; an
estimate that fails any row does not ship.

\paragraph{E-values.} The E-value \citep{vanderweele2017} is the minimum strength of
association, on the risk ratio scale, that an unmeasured confounder would need with \emph{both}
treatment and outcome, above and beyond $Z$, to explain away the observed association. For an
observed risk ratio $\mathrm{RR} \ge 1$,
\begin{equation}\label{eq:causal-evalue}
\mathrm{E} = \mathrm{RR} + \sqrt{\mathrm{RR}\,(\mathrm{RR} - 1)} .
\end{equation}
Our outcome is continuous, so we standardize: with effect $d = \hat\tau / \hat\sigma_Y$
(cluster-bootstrap CI carried through), the conventional conversion $\mathrm{RR} \approx
\exp(0.91\, d)$ feeds \cref{eq:causal-evalue}, and the E-value of the CI limit closer to the
null is reported alongside that of the point estimate. An E-value of 1.8 reads: a latent-quality
confounder must nearly double both the odds of raising and of high growth, conditional on the
full fundamentals history, to explain the estimate away --- a statement a domain expert can
assess directly.

\paragraph{Placebo treatment.} Define $A^{\mathrm{fut}}_{i,y} = \ind{\text{round in year } y +
3}$ and estimate its ``effect'' on $Y_{i,y}$ (which is fully realized before the placebo round)
with the identical AIPW pipeline. Under \cref{ass:causal-ignorability} the estimate must be
null: a future round cannot cause past growth, so any nonzero adjusted estimate measures exactly
the selection-on-trends that $Z$ failed to absorb.

\paragraph{Negative-control outcome.} Symmetrically, estimate the ``effect'' of $A_{i,y}$ on
\emph{pre-treatment} growth (years $y-2$ to $y$, excluded from $Z$ for this test only). Nonzero
means the treated were already on a different trajectory in a way $Z$ does not capture.

\begin{table}[htbp]
\centering
\caption{Mandatory pre-shipping checklist for any causal estimate.}
\label{chca:tab:checklist}
\small
\begin{tabular}{@{}lll@{}}
\toprule
Check & Pass criterion & On failure \\
\midrule
Overlap plot & common support, trimmed & redefine estimand (ATT, \\
 & fraction $< 5\%$ & trimmed population) \\
Placebo (future round) & CI covers 0 & enrich $Z$ (longer history, \\
 & & trend features); else do not ship \\
Negative control (pre-growth) & CI covers 0 & same as placebo \\
E-value & $>$ plausible latent-quality & report as bounded, not \\
 & association given rich $X$ & as an effect \\
Variance agreement & IF vs bootstrap within 20\% & investigate
nuisance \\
 & & instability \\
Synthetic recovery & \cref{chca:tab:synth} ranking & pipeline bug; fix
before \\
 & reproduced & touching real data \\
\bottomrule
\end{tabular}
\end{table}

\section{Synthetic validation protocol}\label{chca:sec:synth}

The synthetic world of \cref{ch02:sec:synthetic} generates firms with a known latent quality
driving both funding intensity and fundamentals (the regime-B generator makes the
quality--funding link explicit), so a true $\tau$ can be injected by construction and
confounding by quality is present by design. The causal layer is accepted only if it passes the
following protocol there; the results are to be reported in \cref{ch:results-synth}, under the
transfer convention fixed there --- synthetic results set priors and vetoes, never substitutes
for real-data evaluation.

\begin{enumerate}
\item \emph{Generate} an EU-style panel: fundamentals as noisy, lagged functions of latent
quality $q_i$ and macro; treatment assignment with probability increasing in $q_i$ and in the
observable dossier; outcome with a known additive effect $\tau^*$ of treatment.
\item \emph{Estimate} with four pipelines: naive difference in means; g-formula (T-learner);
IPW; cross-fitted AIPW --- each with firm-cluster bootstrap CIs and E-values.
\item \emph{Verify} the qualitative ranking of \cref{chca:tab:synth}: the naive contrast
overstates $\tau^*$ badly (quality confounding); all adjusted estimators land near $\tau^*$ when
the generated $X$ is rich; ablating $X$ columns (hiding the dossier, i.e.\ simulating the opaque
regime of \cref{ch:proxies}) re-opens the gap, with AIPW degrading most gracefully.
\item \emph{Report} empirical CI coverage of $\tau^*$ across Monte Carlo replications and check
that the naive estimator's E-value correctly flags fragility while the adjusted estimators'
E-values are informative about the injected residual confounding.
\end{enumerate}

\begin{table}[htbp]
\centering
\caption{Expected qualitative behaviour on the synthetic testbed
(\cref{ch02:sec:synthetic}). ``$\approx$'' means CI covers $\tau^*$ at
nominal rate.}
\label{chca:tab:synth}
\small
\begin{tabular}{@{}lll@{}}
\toprule
Estimator & Rich $X$ (transparent) & Impoverished $X$ (opaque) \\
\midrule
Naive difference & $\gg \tau^*$ & $\gg \tau^*$ \\
g-formula (XGB) & $\approx \tau^*$ & biased toward naive \\
IPW (trimmed) & $\approx \tau^*$, wider CI & biased, unstable \\
AIPW cross-fitted & $\approx \tau^*$, correct coverage & residual bias,
flagged \\
 & & by placebo and E-value \\
\bottomrule
\end{tabular}
\end{table}

This protocol makes the chapter testable: the code that must reproduce \cref{chca:tab:synth} is
byte-identical to the code run on the real panel, so a pipeline that cannot recover a known
effect is caught before it is believed.

\section{Limits}\label{chca:sec:limits}

Two limits are structural. First, \emph{selection on unobservables is not solved} --- it is
attenuated by a rich $X$, bounded by E-values, and probed by placebo and negative-control
tests, but a confounder invisible to three years of audited filings and strong enough to
survive the checklist remains logically possible. The claim is calibrated accordingly: adjusted
estimates with reported sensitivity, not oracle effects. Second, \emph{external validity across
macro regimes}: $M$ enters both selection and outcome, all estimates average over the funding
cycles in sample, and a CATE-by-era analysis is data-hungry because treated firms are scarce
within any single era. An effect estimated across 2015--2021 need not
transport to a tight-money regime, for the same no-extrapolation reasons as in \cref{ch:eufund}.
The opaque regime sharpens both limits at once: with no fundamentals panel there is no rich $X$,
and the honest response is to change the estimand and the observation model --- the programme of
\cref{ch:proxies}, completed by the point-process machinery of the following part.

\begin{decision}{D-EU.2 --- AIPW behind a falsification gate}
The committed causal estimator for the transparent regime is cross-fitted AIPW
\cref{eq:causal-aipw} with firm-level folds, propensity trimming to $[0.02, 0.98]$, and both
cluster-robust influence-function and cluster-bootstrap intervals (the wider is reported);
g-formula and IPW are retained as diagnostics only. No causal estimate ships without clearing
every row of \cref{chca:tab:checklist} --- overlap, placebo, negative control, E-value, variance
agreement, and synthetic recovery. Rationale: double robustness (\cref{prop:causal-dr}) and the
rate condition of \cref{rem:causal-rates} are the only route to $\sqrt{n}$ inference with
boosted-tree nuisances, and the checklist converts the unverifiable
\cref{ass:causal-ignorability} into auditable failure modes. Reversal trigger: if the placebo or
negative-control tests fail persistently even after enriching $Z$ with longer histories and
trend features, the transparent-regime causal claim is demoted to the bounded/associational
reporting of \cref{ch:eufund} and this box is rewritten to say so.
\end{decision}

\begin{codehooks}
The pipeline of \cref{sec:causal-estimators} is a planned module \modrepo{eu/causal.py}:
firm-level fold assignment, out-of-fold nuisance fitting (reusing the boosted stack of
\cref{ch:eufund} with tuning frozen), trimming with recorded fractions, AIPW scores, and the two
variance routes of \cref{sec:causal-inference}. The checklist of \cref{chca:tab:checklist} is an
artifact contract: the module emits the overlap histogram, placebo and negative-control
estimates, E-values, and variance-agreement ratio as machine-readable outputs alongside
$\hat\tau$, in the \code{results/} conventions of \cref{prot:ch07:temporal}(d). The synthetic
protocol of \cref{chca:sec:synth} builds its confounded panel from \code{synthetic/firms.py} and
\code{synthetic/regime\_b.py} (latent quality already drives funding intensity there) and
belongs in the test suite, not in a notebook: recovery of $\tau^*$ is a regression test that
gates changes to the causal module.
\end{codehooks}
